\section*{الفصل \ref{ch:b2:flow-balances} --- موازنتا كمية الحركة والطاقة في الجريانات}

\begin{solution}{exo:b2:flow-balances:1}
$s = \qty{7.1}{cm^2}$، و $v = 0.025/7.1 \times 10^{-4} = \qty{35}{m/s}$؛ و $F = D_mv = 25
\times 35 = \qty{880}{N}$ على الفوهة وعلى الجدار؛ وأما الجدار المتراجع: $\rho s(v -
u)^2 = 1000 \times 7.1 \times 10^{-4} \times 25.4^2 = \qty{460}{N}$.
\end{solution}

\begin{solution}{exo:b2:flow-balances:2}
$F = 250 \times 3000 = \qty{750}{kN}$؛ و $a_0 = 750/60 - 9.81 = \qty{2.7}{m/s^2}$؛ وعند
\qty{20}{t}: $37.5 - 9.8 = \qty{28}{m/s^2}$. و $\Delta v = 3000\ln3 = \qty{3.3}{km/s}$ في
الفضاء الحرّ؛ ناقصًا $g\tau = \qty{1.6}{km/s}$: أي \qty{1.7}{km/s}.
\end{solution}

\begin{solution}{exo:b2:flow-balances:3}
$S = \qty{177}{cm^2}$: فيكون $PS = \qty{7.1}{kN}$، و $\rho Sv^2 = \qty{71}{N}$. وكل مركّبة
من القوة على الكوع هي $PS + \rho Sv^2 = \qty{7.1}{kN}$: أي \qty{10}{kN} على امتداد
منصّف الزاوية نحو الخارج، و $99\%$ منها ضغط. وبمخرج مفتوح (والضغط النسبي
معدوم هناك): $7.1$ kN على امتداد جهة الدخول و \qty{71}{N} على امتداد
المخرج: أي \qty{7.1}{kN}، على محور الدخول تقريبًا.
\end{solution}

\begin{solution}{exo:b2:flow-balances:4}
$F = 80 \times 550 = \qty{44}{kN}$. وفي الطيران: $F = 80(550 - 240) = \qty{25}{kN}$، و $Fv =
\qty{6.0}{MW}$؛ و $\mathcal P_{\text{حركية}} = \tfrac12 \times 80 \times (550^2 - 240^2) =
\qty{9.8}{MW}$؛ و $\eta_p = 0.61 = 2 \times 240/790$.
\end{solution}

\begin{solution}{exo:b2:flow-balances:5}
$D_m = \qty{500}{kg/s}$. (أ) $F = 1000(60 - u)$ نيوتن، و $\mathcal P = 1000u(60 - u)$؛ و $u =
\qty{30}{m/s}$، و $\omega = 37.5$ rad/s $= \qty{360}{rpm}$. (ب) \qty{900}{kW} $=
\tfrac12D_mv^2$: أي $100\%$. (ج) $F = D_m(v - u)(1 + \cos15^\circ) = \qty{29.5}{kN}$، و $\mathcal P
= \qty{885}{kW}$، و $98\%$. (د) $\Gamma = FR = \qty{24}{kN.m}$؛ وعند التوقف تتضاعف
القوة ($\Gamma = \qty{48}{kN.m}$)؛ وعند $u = v$ تنعدم القوة والعزم.
\end{solution}

\begin{solution}{exo:b2:flow-balances:6}
(أ) $s = \qty{7.1}{mm^2}$ لكل فوهة، و \qty{0.10}{L/s} لكل منهما: ومنه $w = \qty{14}{m/s}$. (ب)
$\Gamma_0 = D_mRw = 0.2 \times 0.15 \times 14 = \qty{0.42}{N.m}$. (ج) $D_mR(w - R\omega) =
5 \times 10^{-3}$: ومنه $w - R\omega = \qty{0.17}{m/s}$، و $\omega = \qty{93}{rad/s} \approx \qty{890}{rpm}$؛
ودون احتكاك $R\omega = w$: أي \qty{900}{rpm}. (د) \qty{0.17}{m/s} إلى الخلف، ثم
صفر: فيسقط الماء شاقوليًّا.
\end{solution}

\begin{solution}{exo:b2:flow-balances:7}
(أ) $D_m = \rho Av_i$، \omterm{thm:b2:flow-balances:momentum}{وتدفق كمية الحركة} $D_m\cdot2v_i$: ومنه $T = 2\rho Av_i^2$؛ و $A = \qty{78.5}{m^2}$،
و $T = \qty{19.6}{kN}$: أي $v_i = \qty{10}{m/s}$. (ب) $\mathcal P = \tfrac12D_m(2v_i)^2 = Tv_i =
\qty{200}{kW}$. (ج) $\mathcal P = T^{3/2}/\sqrt{2\rho A}$: فمضاعفة المساحة
تقسم الاستطاعة المحرَّضة على $\sqrt2$ --- إذ يحرّك دوّار أكبر هواءً أكثر
أبطأ. (د) $A = \qty{0.20}{m^2}$، و $v_i = \sqrt{9.81/0.47} = \qty{4.6}{m/s}$، و $\mathcal P
= \qty{45}{W}$: أي \qty{45}{W/kg} مقابل \qty{100}{W/kg} --- فحمولة القرص $T/A$
أدنى.
\end{solution}

\begin{solution}{exo:b2:flow-balances:8}
(أ) $v_2 = \qty{1.0}{m/s}$؛ وببرنولي $\tfrac12\rho(16 - 1) = \qty{7.5}{kPa}$؛ والفعلي
$\rho v_2(v_1 - v_2) = \qty{3.0}{kPa}$؛ والضياع \qty{4.5}{kPa} $= \qty{0.46}{m}$. (ب) $6$ kPa
مستعادة، و \qty{1.5}{kPa} ضائعة. (ج) $\tfrac12\rho v_1^2 = \qty{8}{kPa}$، أي \qty{0.82}{m}. (د)
$D_V = \qty{7.9}{L/s}$: أي \qty{35}{W} و \qty{63}{W}.
\end{solution}

\begin{solution}{exo:b2:flow-balances:9}
(أ) $v = 0.01/2.83 \times 10^{-3} = \qty{3.5}{m/s}$، و $v^2/2g = \qty{0.64}{m}$؛ والاحتكاك
$0.025 \times 1000 \times 0.64 = \qty{16}{m}$، والوصلات \qty{1.3}{m}: أي \qty{17.3}{m}. (ب)
$H = 30 + 17.3 + 0.64 = \qty{48}{m}$؛ و $\rho gD_VH = \qty{4.7}{kW}$؛ و \qty{7.2}{kW} عند
العمود. (ج) المدخل (بلا ضياعات سحب): $P_0 - \rho g(5 + 0.64) = \qty{45}{kPa}$؛
والمخرج $45 + 470 = \qty{515}{kPa}$. (د) $P_{\text{المدخل}} \ge \qty{2.3}{kPa}$: ومنه $h \le (100 -
2.3)/9.81 - 0.64 = \qty{9.3}{m}$ (وأقل مع ضياعات السحب؛ و $7$--\qty{8}{m}
عمليًّا).
\end{solution}

\begin{solution}{exo:b2:flow-balances:10}
(أ) $v = \sqrt{2 \times 4 \times 10^5/1000} = \qty{28}{m/s}$؛ و $s = \qty{3.8}{cm^2}$، و $D_m =
\qty{10.8}{kg/s}$، و $F = D_mv = 2\Delta P\,s = \qty{300}{N}$؛ و $a = 300/1.1 - 9.8 =
\qty{260}{m/s^2}$. (ب) والهواء عند \qty{1.5}{L}: $P = \qty{3.3}{bar}$، و $v = \sqrt{2 \times 2.3 \times
10^5/1000} = \qty{22}{m/s}$. (ج) وسرعة الخروج الوسطى $\approx \qty{22}{m/s}$: فالاندفاع $\approx 10^{-3}
/(3.8 \times 10^{-4} \times 22) = \qty{0.12}{s}$؛ وبمعادلة الصاروخ مع $v_e \approx \qty{22}{m/s}$:
$\Delta v \approx 22\ln(1.1/0.1) = \qty{50}{m/s}$. (د) وينفث الهواء الباقي عند \qty{2.5}{bar}
فيضيف قليلًا؛ وكثرة الماء تترك هواءً قليلًا فينهار
الضغط باكرًا، وقلّته تترك كتلة قليلة تُرمى ---
والثلث من الحجم هو الأفضل.
\end{solution}

\begin{solution}{exo:b2:flow-balances:11}
(أ) $\Gamma = D_mr_1v_{\theta1} = 3 \times 10^4 \times 20 = \qty{600}{kN.m}$، و $\omega = 15.7$ rad/s،
و $\mathcal P = \qty{9.4}{MW}$. (ب) $H = \mathcal P/\rho gD_V = \qty{32}{m}$. (ج) $\Gamma = 3 \times
10^4(20 - 1.5) = \qty{555}{kN.m}$، و $\mathcal P = \qty{11.6}{MW}$ (والعلو \qty{39}{m})؛ وتحمل
دوّامة الخروج $v_{\theta2}^2/2g = \qty{0.46}{m}$، أي نحو $1\%$. (د) ونفث بلتون
عند الضغط الجوي: فيحتاج إلى علو كبير ليكون سريعًا، وتدفقه
محدود بالفوهة؛ وأما دوّار فرانسيس فمملوء بالماء تحت
ضغط ويمرّر تدفقات كبيرة عند علوّات معتدلة.
\end{solution}

\begin{solution}{exo:b2:flow-balances:12}
(أ) $v_1h_1 = v_2h_2 = q$. (ب) القوة الهيدروستاتيكية في واحدة العرض على
مقطع $\int_0^h\rho g(h - z)\dd z = \tfrac12\rho gh^2$؛ وكمية الحركة: $\tfrac12\rho g(h_1^2
- h_2^2) = \rho q(v_2 - v_1)$؛ ثم نقسم على $\rho$ ونستعمل $q = v_1h_1 = v_2h_2$.
(ج) ومع $v_2 = v_1h_1/h_2$ و $h_1 \ne h_2$: $\tfrac12g(h_1 + h_2)h_2 = v_1^2h_1$،
أي $(h_2/h_1)^2 + h_2/h_1 - 2\mathrm{Fr}_1^2 = 0$؛ ويقتضي $h_2 > h_1$ أن يكون $\mathrm{Fr}_1 > 1$.
(د) $\Delta H = (h_1 + v_1^2/2g) - (h_2 + v_2^2/2g) = (h_2 - h_1)^3/4h_1h_2$.
و $\mathrm{Fr}_1 = 2.86$، و $h_2/h_1 = 3.57$، و $h_2 = \qty{0.71}{m}$، و $v_2 = \qty{1.1}{m/s}$،
و $\Delta H = 0.514^3/0.57 = \qty{0.24}{m}$؛ و $\rho gq\,\Delta H = \qty{1.9}{kW}$ في المتر
من العرض.
\end{solution}

\begin{solution}{pb:b2:flow-balances:1}
\textbf{1.} $F = (D_a + D_f)v_e = 51 \times 600 = \qty{30.6}{kN}$.

\textbf{2.} $F = 51 \times 600 - 50 \times 250 = \qty{18.1}{kN}$؛ و $Fv = \qty{4.5}{MW}$.

\textbf{3.} $\mathcal P_{\text{حركية}} = \tfrac12 \times 51 \times 600^2 - \tfrac12 \times 50 \times 250^2 =
\qty{7.6}{MW}$؛ و $\eta_p = 0.59$. ومع $D_f = 0$: $\eta_p = 2v(v_e - v)/(v_e^2 - v^2)
= 2v/(v_e + v) = 500/850$.

\textbf{4.} الحرارة \qty{43}{MW}: فالحراري $7.6/43 = 18\%$؛ والكلي $4.5/43 = 10\%$.

\textbf{5.} $\tfrac12 \times 500(v_e'^2 - 250^2) = 7.6 \times 10^6$: ومنه $v_e' = \qty{305}{m/s}$؛
و $F = 500 \times 55 = \qty{27.5}{kN}$؛ و $\eta_p = 500/555 = 0.90$. فتحريك هواء أكثر
أبطأ يعطي دفعًا أكبر بالاستطاعة نفسها: ومن هنا المراوح.

\textbf{6.} لأنه يحتاج إلى الهواء الخارجي مائعًا عاملًا ومؤكسدًا معًا؛
وأما الصاروخ فيحمل الاثنين.

\textbf{7.} يدخل الهواء بالسرعة \qty{70}{m/s} (إلى الخلف في مرجع المحرك)
ويغادر بمركّبة أمامية $150\cos45^\circ = \qty{106}{m/s}$: فكمية الحركة
الممنوحة للهواء إلى الأمام، أي $500(106 + 70) = \qty{88}{kN}$ من
الكبح؛ والجريان نفسه مقذوفًا إلى الخلف كان ليعطي $500(150 - 70) =
\qty{40}{kN}$ من الدفع.

\textbf{8.} على امتداد $\dd t$ تعطي الجملة (الصاروخ $m$ مع الغاز $D_m\dd t$ المقذوف عند
$v - v_e$): $(m - D_m\dd t)(v + \dd v) + D_m\dd t(v - v_e) - mv = -mg\,\dd t$،
ومنه $m\,\dd v/\dd t = D_mv_e - mg$.

\textbf{9.} $D_m = \qty{1333}{kg/s}$؛ و $F = \qty{4.0}{MN}$؛ و $a = 13.3 - 9.8 =
\qty{3.5}{m/s^2}$ عند الإقلاع، و $40 - 9.8 = \qty{30}{m/s^2}$ عند نهاية الاندفاع.

\textbf{10.} $v = v_e\ln(m_0/m) - gt$: أي $3000\ln3 - 9.81 \times 150 = 3296 - 1472 =
\qty{1.8}{km/s}$؛ وضياع الثقالة \qty{1.5}{km/s}.

\textbf{11.} المرحلة 1: $3000\ln(300/200) = \qty{1.2}{km/s}$؛ ثم إلقاء \qty{20}{t}؛ والمرحلة
2: $3000\ln(180/80) = \qty{2.4}{km/s}$: أي \qty{3.6}{km/s} مقابل \qty{3.3}{km/s}
للمرحلة الواحدة (ومع الثقالة تتسع الفجوة).

\textbf{12.} $v_e = 7800/\ln3 = \qty{7.1}{km/s}$: ولا يبلغها أي وقود دافع
كيميائي (فالهيدروجين مع الأكسجين يعطي \qty{4.5}{km/s})؛ فالتقسيم مراحل واجب.

\textbf{13.} يُنزع \omterm{thm:b2:flow-balances:momentum}{تدفق كمية الحركة} الشاقولي للنفث $D_mv = \qty{4.0}{MN}$
ويُنشأ مثله أفقيًّا: أي قوة \qty{4}{MN} نحو الأسفل
مع \qty{4}{MN} جانبيًّا، و \qty{5.7}{MN} في المجموع.

\textbf{14.} تتضمن الموازنة على المحرك الحدَّ $(P_e - P_0)S_e$ على
مقطع المخرج؛ وبهبوط $P_0$ مع الارتفاع ينمو هذا الحدّ: فيرتفع الدفع
بنحو $10$--$15\%$ في الخلاء.

\textbf{15.} $D_m = \qty{3e4}{kg/s}$ لكل عجلة: $F = D_m(v - u)(1 + \cos15^\circ) =
1.97 \times 3 \times 10^4(62.6 - u)$؛ و $\mathcal P = Fu$، أعظمية عند $u = v/2 = \qty{31.3}{m/s}$.

\textbf{16.} $\omega = u/R = 20.9$ rad/s $= \qty{199}{rpm}$؛ و $F = \qty{1.85}{MN}$، و $\Gamma = FR
= \qty{2.8}{MN.m}$.

\textbf{17.} $f = pn/60$: ومنه $p = 3000/199 \approx 15$ زوج أقطاب، و $n = \qty{200}{rpm}$.

\textbf{18.} $\mathcal P = Fu = \qty{58}{MW}$؛ واستطاعة النفث $\tfrac12D_mv^2 = \qty{59}{MW}$:
أي $98\%$؛ و $2 \times 58 \times 0.95 = \qty{110}{MW}$.

\textbf{19.} الداخل: $Rv = \qty{94}{m^2/s}$ في الكيلوغرام. والخارج: السرعة السمتية المطلقة
$u - (v - u)\cos15^\circ = \qty{1.1}{m/s}$، أي \qty{1.6}{m^2/s}. ومنه $\Gamma = D_m(94 -
1.6) = \qty{2.8}{MN.m}$ --- كما في السؤال 16.

\textbf{20.} حين $u \to v$ تنعدم السرعة النسبية والقوة والعزم:
فلا شيء يكبح العجلة، فتفرط في السرعة؛ ويرمي الحارف
النفث بعيدًا عن الأكواب في أقل من ثانية بينما ينغلق صمام الإبرة
ببطء يكفي لتجنّب الصدمة المائية.

\textbf{21.} $v = \qty{5.7}{m/s}$، و $v^2/2g = \qty{1.6}{m}$؛ والضياع $0.015 \times 133 \times 1.6 =
\qty{3.3}{m}$؛ وعلو المضخة $200 + 3.3 + 1.6 = \qty{205}{m}$.

\textbf{22.} $\rho gD_VH = 9810 \times 40 \times 205 = \qty{80}{MW}$؛ و \qty{91}{MW} كهربائية.

\textbf{23.} $P = P_0 + \rho g(200 + 3.3) + \tfrac12\rho v^2 \approx \qty{21}{bar}$ مطلقًا،
مقابل \qty{20.3}{bar} في نظام العنفة: فالضياعات تُضاف الآن إلى
العلو السكوني بدل أن تُطرح منه.

\textbf{24.} $V = 40 \times 28800 = \qty{1.15e6}{m^3}$، و $E = \rho gV \times 200 = \qty{2.3e12}{J}
= \qty{630}{MWh}$؛ والمستهلَك $91 \times 8 = \qty{730}{MWh}$؛ والمستعاد $630 \times 0.9 \times
0.95 = \qty{540}{MWh}$: فمردود الدورة الكاملة $74\%$.

\textbf{25.} كمية الحركة: دفع المحرك النفاث ودفع الصاروخ، وقوة
بلتون. والعزم الحركي: عزم العنفة (أويلر). والطاقة:
علو المضخة واستطاعتها، ومردود التخزين. والكتلة: كل تدفق،
وفقد الصاروخ لكتلته.
\end{solution}
